\documentclass[../../../include/open-logic-section]{subfiles}

\begin{document}

\olfileid[hi]{sth}{choice}{woproblem}
\olsection{सुव्यवस्था समस्या}

स्पष्टतः बहुत कुछ इस पर निर्भर है कि हम सुव्यवस्था को स्वीकार करते हैं या
नहीं। किंतु इस सिद्धांत की चर्चा सामान्यतः उसके तुल्य सिद्धांत, चयन
अभिगृहीत, पर केंद्रित रही है। इसलिए अब हम उस ओर ध्यान देंगे (और तुल्यता
सिद्ध करेंगे)।

\citeyear{Cantor1883} में कैंटर ने सुव्यवस्था अभिगृहीत का समर्थन करते हुए
उसे ``विचार का ऐसा नियम जो मुझे मूलभूत, अपने परिणामों में समृद्ध और अपनी
सामान्य वैधता के लिए विशेष रूप से उल्लेखनीय प्रतीत होता है'' कहा
(\citeauthor{Potter2004} \citeyear[p.~243]{Potter2004} में उद्धृत)। किंतु
अंततः कैंटर आश्वस्त हो गए कि इस ``अभिगृहीत'' को प्रमाण की आवश्यकता है।
गणितीय समुदाय भी ऐसा ही मानने लगा।

इस समस्या को ज़रमेलो ने \citeyear{Zermelo1904} में ``हल'' किया। उनके
समाधान की व्याख्या के लिए हमें कुछ परिभाषाएँ चाहिए।

\begin{defn}
कोई फलन $f$ एक \emph{चयन फलन} है तभी जब सभी $x\in\dom{f}$ के लिए
$f(x)\in x$। हम कहते हैं कि $f$, $A$ \emph{के लिए चयन फलन} है तभी जब $f$
ऐसा चयन फलन है जिसके लिए $\dom{f}=A\setminus\{\emptyset\}$।
\end{defn}

सहज रूप से $A$ के लिए चयन फलन प्रत्येक (अरिक्त) समुच्चय $x\in A$ से एक
विशिष्ट अवयव $f(x)$ \emph{चुनता} है। तब चयन अभिगृहीत है:

\begin{axiom}[चयन]
	प्रत्येक समुच्चय का कोई चयन फलन है।
\end{axiom}

ज़रमेलो ने दिखाया कि चयन से सुव्यवस्था निष्पन्न होती है और इसका विलोम भी:

\begin{thm}[$\ZF$ में]\ollabel{thmwochoice}
सुव्यवस्था और चयन तुल्य हैं।
\end{thm}

\begin{proof}
\emph{बाएँ-से-दाएँ।} $A$ को समुच्चयों का समुच्चय मानें। तब सम्मिलन
अभिगृहीत से $\bigcup A$ का अस्तित्व है, अतः सुव्यवस्था से कोई $<$ है जो
$\bigcup A$ को सुव्यवस्थित करता है। अब
$f(x)=\text{$x$ का $<$-लघुतम सदस्य}$ रखें। यह $A$ के लिए चयन फलन है।

\emph{दाएँ-से-बाएँ।} $A$ को नियत करें। चयन से
$\Pow{A}\setminus\{\emptyset\}$ के लिए कोई चयन फलन $f$ है। परिमितेतर
पुनरावर्तन का उपयोग करके एक फलन परिभाषित करें:
\begin{align*}
	g(0) &= f(A)\\
	g(\alpha) &= 
		\begin{cases}
			\text{रुकें!{}} &\text{यदि }A = \funimage{g}{\alpha}\\
			f(A \setminus \funimage{g}{\alpha}) & \text{अन्यथा}\\	
		\end{cases}
\end{align*}
``रुकें!'' का संकेत केवल उस परिभाषा का संक्षिप्त रूप है जो अन्यथा अधिक
लंबी होती। अर्थात् पहली बार जब $A=\funimage{g}{\alpha}$ हो, तो सभी
$\delta\leq\alpha$ के लिए $g(\delta)=A$ रखें। अभी आरंभ में हम केवल इतना
निश्चित हो सकते हैं कि इससे कोई \emph{पद} परिभाषित होता है
(\olref[sth][spine][recursion]{transrecursionschema} के बाद की टिप्पणियाँ
देखें); किंतु हम दिखाएँगे कि हमारे पास वास्तव में एक फलन है।

चूँकि $f$ चयन फलन है, प्रत्येक $\alpha$ के लिए (जब वह परिभाषित हो) हमारे
पास $g(\alpha)=f(A\setminus\funimage{g}{\alpha})\in
A\setminus\funimage{g}{\alpha}$; अर्थात्
$g(\alpha)\notin\funimage{g}{\alpha}$। अतः यदि
$g(\alpha)=g(\beta)$, तो $g(\beta)\notin\funimage{g}{\alpha}$, अर्थात्
$\beta\notin\alpha$, और इसी प्रकार $\alpha\notin\beta$। अतः त्रिविभाजन से
$\alpha=\beta$। इसलिए $g$ !!{injective} है।

इसके बाद ध्यान दें कि हम सच में रुकते हैं!{}, अर्थात् कोई (लघुतम) क्रमसूचक
$\alpha$ है ऐसा कि $A=g[\alpha]$। अन्यथा मानें; तब $g$ के !!{injective}
होने से प्रत्येक क्रमसूचक $\alpha$ के लिए
$\cardless{\alpha}{\Pow{A}\setminus\{\emptyset\}}$ होता, जो
\olref[hartogs]{HartogsLemma} के विरुद्ध है। अतः $\ran{g}=A$ भी।

इन तथ्यों को मिलाकर $g$ किसी क्रमसूचक से $A$ तक !!a{bijection} है। अब $g$
का उपयोग $A$ को सुव्यवस्थित करने में किया जा सकता है।
\end{proof}

अतः सुव्यवस्था और चयन साथ ही टिकते या गिरते हैं। किंतु प्रश्न शेष है: वे
टिकते हैं या गिरते हैं?

\end{document}
